\documentclass[11pt]{article}
\usepackage[margin=1in]{geometry}
\usepackage{amsmath,amssymb,amsthm}
\usepackage{enumitem}

\newtheorem{theorem}{Theorem}
\newtheorem{lemma}[theorem]{Lemma}
\newtheorem{proposition}[theorem]{Proposition}
\newtheorem{corollary}[theorem]{Corollary}
\newtheorem{claim}[theorem]{Claim}

\theoremstyle{definition}
\newtheorem{definition}[theorem]{Definition}
\newtheorem{question}{Question}
\newtheorem{remark}[theorem]{Remark}

\title{Analysis of the Corridor Condition:\\Stability, Blow-up Constraints, and Natural Satisfaction}
\author{Technical Analysis}
\date{}

\begin{document}
\maketitle

\section{The Corridor Condition}

From the paper, the ``corridor condition'' is the requirement that under concentration ($\rho \geq 1 - \varepsilon_*$), the characteristic frequency $\lambda = \sqrt{D/E}$ satisfies the dichotomy bounds derived in the Closing Lemma.

\begin{definition}[Corridor Condition]
A solution $u(t)$ satisfies the \textbf{corridor condition} at time $t$ if one of the following holds:
\begin{enumerate}[label=(\roman*)}
\item \textbf{Non-concentrated:} $\rho(t) < 1 - \varepsilon_*$
\item \textbf{Concentrated, low frequency:} $\rho(t) \geq 1 - \varepsilon_*$ and $\lambda(t) < \lambda_{crit} := c_1/(4\nu)$
\item \textbf{Concentrated, high frequency:} $\rho(t) \geq 1 - \varepsilon_*$ and $\lambda(t) \geq \lambda_{crit}$
\end{enumerate}
where $E = \|u\|_{\dot{H}^{1/2}}^2$, $D = \|u\|_{\dot{H}^{3/2}}^2$, $\lambda = \sqrt{D/E}$, and $\rho = E_{max}/E$.
\end{definition}

The key claim is that ALL solutions satisfy the corridor condition (by definition---it's exhaustive), but each case leads to a different consequence:
\begin{itemize}
\item Case (i): Standard regime, no special analysis needed
\item Case (ii): $d\rho/dt < 0$, concentration breaks
\item Case (iii): $dE/dt < 0$, energy decreases
\end{itemize}

%==============================================================================
\section{Question 1: Forward Stability of the Corridor Condition}
%==============================================================================

\begin{question}
Is the corridor condition forward-stable under NS evolution? Can it fail instantaneously, or must it degrade continuously?
\end{question}

\subsection{Analysis}

First, note that the corridor condition as defined is \textbf{always satisfied}---it's a partition of the state space, not a constraint. The real question is:

\textbf{Can a solution jump discontinuously between cases (i), (ii), (iii)?}

The relevant quantities are:
\begin{itemize}
\item $\rho(t) = E_{max}(t)/E(t)$ (concentration ratio)
\item $\lambda(t) = \sqrt{D(t)/E(t)}$ (characteristic frequency)
\end{itemize}

\begin{proposition}[Continuity of State Variables]
For smooth NS solutions, $\rho(t)$ and $\lambda(t)$ are continuous functions of time.
\end{proposition}

\begin{proof}
$E(t) = \|u(t)\|_{\dot{H}^{1/2}}^2$ and $D(t) = \|u(t)\|_{\dot{H}^{3/2}}^2$ are continuous for smooth solutions (they're integrals of continuous integrands).

For $E_{max}(t) = \max_j E_j(t)$: each $E_j(t)$ is continuous, and the maximum of finitely many continuous functions is continuous.

Therefore $\rho = E_{max}/E$ and $\lambda = \sqrt{D/E}$ are continuous (assuming $E > 0$, which holds for non-trivial solutions).
\end{proof}

\begin{corollary}[No Instantaneous Failure]
A solution cannot jump instantaneously between cases. Transitions occur continuously.
\end{corollary}

\subsection{Propagation Result}

\begin{theorem}[Short-Time Stability of Case Classification]\label{thm:short-time}
Let $u(t)$ be a smooth NS solution.

\textbf{(a)} If $\rho(t_0) < 1 - \varepsilon_* - \delta$ for some $\delta > 0$, then there exists $\tau > 0$ such that $\rho(t) < 1 - \varepsilon_*$ for all $t \in [t_0, t_0 + \tau]$.

\textbf{(b)} If $\rho(t_0) \geq 1 - \varepsilon_*$ and $\lambda(t_0) < \lambda_{crit} - \delta$, then there exists $\tau > 0$ such that $\lambda(t) < \lambda_{crit}$ for all $t \in [t_0, t_0 + \tau]$.

\textbf{(c)} If $\rho(t_0) \geq 1 - \varepsilon_*$ and $\lambda(t_0) > \lambda_{crit} + \delta$, then there exists $\tau > 0$ such that $\lambda(t) > \lambda_{crit}$ for all $t \in [t_0, t_0 + \tau]$.

Moreover, $\tau$ can be estimated in terms of the rates $|d\rho/dt|$ and $|d\lambda/dt|$.
\end{theorem}

\begin{proof}
All three statements follow from continuity. For quantitative estimates:

\textbf{(a)} We have $|d\rho/dt| \leq C(\rho, E, D)$ bounded for smooth solutions. The time to cross from $\rho_0 < 1 - \varepsilon_* - \delta$ to $1 - \varepsilon_*$ is at least $\tau \geq \delta / \sup|d\rho/dt|$.

From the paper's analysis:
\[
\frac{d\rho}{dt} = \frac{1}{E}\left(\frac{dE_{max}}{dt} - \rho\frac{dE}{dt}\right)
\]

Both $dE_{max}/dt$ and $dE/dt$ involve trilinear forms bounded by $O(ED^{1/2})$.

So $|d\rho/dt| \leq C \cdot D^{1/2} = C\lambda E^{1/2}$, giving:
\[
\tau \geq \frac{\delta}{C\lambda E^{1/2}}
\]

\textbf{(b), (c)} Similar analysis for $\lambda$. We have:
\[
\frac{d\lambda}{dt} = \frac{1}{2\lambda E}(\dot{D} - \lambda^2 \dot{E})
\]

Both $\dot{D}$ and $\dot{E}$ are bounded by trilinear + dissipation terms, giving $|d\lambda/dt| \leq C'(\lambda, E, D)$.
\end{proof}

\subsection{Stability Under Perturbation}

\begin{proposition}[Structural Stability]
The case classification is structurally stable: small perturbations in initial data lead to small perturbations in the time spent in each case.
\end{proposition}

\begin{proof}
NS is locally well-posed in $\dot{H}^{1/2}$ (Koch-Tataru). Solutions depend continuously on initial data for short times. Since $\rho$ and $\lambda$ depend continuously on $u$, the case classification is stable under perturbation.
\end{proof}

%==============================================================================
\section{Question 2: Blow-up Analysis---Which Term Diverges First?}
%==============================================================================

\begin{question}
Assume finite-time blow-up. Which term in the dynamics must diverge first? Does this contradict the corridor condition?
\end{question}

\subsection{Setup}

Assume $u(t)$ blows up at time $T^* < \infty$. By definition of critical blow-up:
\[
\|u(t)\|_{\dot{H}^{1/2}} \to \infty \quad \text{as } t \to T^*
\]

The relevant evolution equations are:
\begin{align}
\frac{dE}{dt} &= -2\mathcal{T} - 2\nu D \label{eq:E-evol} \\
\frac{dD}{dt} &= -2\mathcal{T}_{3/2} - 2\nu F \label{eq:D-evol}
\end{align}
where:
\begin{itemize}
\item $\mathcal{T} = \langle (u \cdot \nabla)u, |\nabla|u \rangle$ (trilinear at $\dot{H}^{1/2}$)
\item $\mathcal{T}_{3/2} = \langle (u \cdot \nabla)u, |\nabla|^3 u \rangle$ (trilinear at $\dot{H}^{3/2}$)
\item $F = \|u\|_{\dot{H}^{5/2}}^2$
\end{itemize}

\subsection{Order of Divergence}

\begin{proposition}[Hierarchy of Blow-up]\label{prop:hierarchy}
If $E(t) \to \infty$ as $t \to T^*$, then:
\[
D(t) \to \infty, \quad F(t) \to \infty, \quad \text{and } |\mathcal{T}(t)| \to \infty
\]
Moreover, $D$ diverges at least as fast as $E$.
\end{proposition}

\begin{proof}
By interpolation:
\[
E = \|u\|_{\dot{H}^{1/2}}^2 \leq \|u\|_{L^2} \|u\|_{\dot{H}^1}
\]

Since $\|u\|_{L^2}$ decreases (energy dissipation), $\|u\|_{\dot{H}^1} \to \infty$ if $E \to \infty$.

Similarly, $D = \|u\|_{\dot{H}^{3/2}}^2 \geq c \|u\|_{\dot{H}^{1/2}}^{4/3} / \|u\|_{L^2}^{1/3}$ by interpolation, so $D \to \infty$.

For the trilinear: blow-up requires $dE/dt > 0$ eventually, so from \eqref{eq:E-evol}:
\[
\mathcal{T} < -\nu D \to -\infty
\]
\end{proof}

\begin{proposition}[Rate Constraints]\label{prop:rates}
For blow-up with $E \to \infty$:
\begin{enumerate}
\item $\mathcal{T} < -\nu D$ (trilinear must be large and negative)
\item $|\mathcal{T}| \leq CED^{1/2}$ (upper bound)
\item Combined: $\nu D < CED^{1/2}$, so $\lambda < CE^{1/2}/\nu$
\end{enumerate}
\end{proposition}

\subsection{The Contradiction}

\begin{theorem}[Blow-up Contradicts Corridor Structure]\label{thm:contradiction}
Assume finite-time blow-up with $E(t) \to \infty$ as $t \to T^*$. Then:

\textbf{(a)} The solution must eventually be in the concentrated regime ($\rho \geq 1 - \varepsilon_*$) for times near $T^*$.

\textbf{(b)} Under concentration with blow-up, we need $\lambda < CE^{1/2}/\nu$ (from Proposition~\ref{prop:rates}).

\textbf{(c)} The Closing Lemma states:
\begin{itemize}
\item If $\lambda < \lambda_{crit}$: $d\rho/dt < 0$ (concentration decreases)
\item If $\lambda \geq \lambda_{crit}$: $dE/dt < 0$ (energy decreases)
\end{itemize}

\textbf{(d)} For large $E$, the condition $\lambda < CE^{1/2}/\nu$ is compatible with $\lambda \geq \lambda_{crit}$.

\textbf{(e)} But $\lambda \geq \lambda_{crit}$ implies $dE/dt < 0$, contradicting blow-up.

\textbf{(f)} And $\lambda < \lambda_{crit}$ implies concentration breaks before blow-up completes.
\end{theorem}

\begin{proof}
\textbf{(a)} By concentration-compactness (Gallagher-Koch-Planchon), minimal blow-up elements concentrate. If blow-up occurs without concentration, the solution would decompose into concentrating profiles, and at least one must blow up---but that one concentrates.

\textbf{(b)} From Proposition~\ref{prop:rates}, blow-up requires $\mathcal{T} < -\nu D$, and $|\mathcal{T}| \leq CED^{1/2}$, giving $\lambda < CE^{1/2}/\nu$.

\textbf{(c)} This is the content of the Closing Lemma (Steps 9-10).

\textbf{(d)} For $E > \nu^2\lambda_{crit}^2/C^2 = c_1^2/(16C^2)$, the range $[\lambda_{crit}, CE^{1/2}/\nu]$ is non-empty.

\textbf{(e)} If $\lambda \in [\lambda_{crit}, CE^{1/2}/\nu]$, then from the Closing Lemma:
\[
\frac{dE}{dt} \leq -\frac{c_1^2\rho E}{8\nu} < 0
\]
This contradicts $dE/dt > 0$ required for blow-up.

\textbf{(f)} If $\lambda < \lambda_{crit}$, then $d\rho/dt < 0$. The concentration ratio decreases at rate:
\[
\frac{d\rho}{dt} \leq -c_*\nu\lambda^2(1-\rho) < 0
\]
Eventually $\rho$ drops below $1 - \varepsilon_*$, exiting concentration.
\end{proof}

\subsection{Which Term Diverges First?}

\begin{corollary}
Under the hypothetical blow-up scenario:
\begin{enumerate}
\item $E$ must grow: $E \to \infty$
\item $D$ must grow: $D \geq cE$ (at least linear in $E$)
\item $\mathcal{T}$ must be large negative: $\mathcal{T} < -\nu D$
\item $\lambda$ must satisfy $\lambda < CE^{1/2}/\nu$ (sublinear in $E^{1/2}$)
\end{enumerate}

The term that ``fails first'' is the trilinear $\mathcal{T}$: it needs to be more negative than $-\nu D$, but this is incompatible with:
\begin{itemize}
\item The Closing Lemma bound (under concentration), which gives $dE/dt < 0$ for high $\lambda$
\item The outflow mechanism (under concentration), which breaks concentration for low $\lambda$
\end{itemize}
\end{corollary}

%==============================================================================
\section{Question 3: Natural Satisfaction of the Corridor Condition}
%==============================================================================

\begin{question}
Is the corridor condition automatically satisfied for any nontrivial subclass of NS solutions?
\end{question}

\subsection{Clarification}

The corridor condition (as a partition) is always satisfied. The meaningful question is:

\textbf{For which solutions does the dichotomy immediately imply no blow-up?}

That is: which solutions are always in Case (i) (non-concentrated) or always in Case (iii) (concentrated + high frequency)?

\subsection{Class 1: Almost-Helical Flows}

\begin{definition}
A flow $u$ is $\varepsilon$-helical if $\|u^-\|_{\dot{H}^{1/2}} \leq \varepsilon\|u\|_{\dot{H}^{1/2}}$ where $u = u^+ + u^-$ is the helicity decomposition.
\end{definition}

\begin{proposition}[Almost-Helical Implies Spread]\label{prop:helical-spread}
If $u$ is $\varepsilon$-helical with $\varepsilon < \varepsilon_0$ (small), then $u$ cannot be angularly concentrated.
\end{proposition}

\begin{proof}
For single-helicity $u^+$, the Fourier constraint is $\hat{u}^+(\xi) = a(\xi) h^+(\hat{\xi})$ where $h^+$ is the positive helicity eigenvector: $h^+(\hat{\xi}) = \frac{1}{\sqrt{2}}(e_\theta + ie_\phi)$ in spherical coordinates.

As $\hat{\xi}$ varies over $S^2$, $h^+(\hat{\xi})$ traces out a circle in the plane perpendicular to $\hat{\xi}$.

For concentration in a cap $\Omega$ centered at $\hat{n}$, we'd need $\hat{u}^+(\xi)$ to have a preferred direction. But $h^+(\hat{\xi})$ rotates as $\hat{\xi}$ moves within $\Omega$---there's no fixed direction.

More precisely: the helicity constraint forces
\[
\int_\Omega |h^+(\hat{\xi}) \cdot \hat{e}|^2 d\hat{\xi} = \frac{1}{2}\text{Area}(\Omega)
\]
for ANY fixed direction $\hat{e}$. This means energy cannot concentrate in direction.

For almost-helical flows ($\varepsilon$ small), this near-isotropy persists.
\end{proof}

\begin{corollary}
Almost-helical flows remain in the non-concentrated regime (Case (i)), where standard bounds apply and blow-up is prevented.
\end{corollary}

\subsection{Class 2: Spectrally Localized Data}

\begin{definition}
Data $u_0$ is spectrally localized if $\text{supp}(\hat{u}_0) \subset \{\xi : |\xi| \in [a, b]\}$ for some $0 < a < b < \infty$.
\end{definition}

\begin{proposition}[Spectral Localization Gives High Frequency]\label{prop:spectral}
If $u_0$ is spectrally localized with $a > 0$, then $\lambda(0) \geq a$ and the solution satisfies $\lambda(t) \geq a \cdot e^{-C\nu t}$ for short times.
\end{proposition}

\begin{proof}
For spectrally localized data with $|\xi| \geq a$:
\[
D = \int |\xi|^3 |\hat{u}|^2 d\xi \geq a^2 \int |\xi| |\hat{u}|^2 d\xi = a^2 E
\]
So $\lambda = \sqrt{D/E} \geq a$.

Under NS evolution, the support spreads (low frequencies are generated), but at a rate controlled by the nonlinearity. Standard estimates give the exponential lower bound.
\end{proof}

\begin{corollary}
For spectrally localized data with $a > \lambda_{crit} = c_1/(4\nu)$:
\begin{itemize}
\item If concentrated: $\lambda \geq a > \lambda_{crit}$, so $dE/dt < 0$ (energy decreases)
\item If not concentrated: standard bounds apply
\end{itemize}
Either way, no blow-up.
\end{corollary}

\subsection{Class 3: Dominant Projected Component}

\begin{definition}
Let $\Pi_\Omega$ be the projection onto a cap $\Omega$. Data $u_0$ has $(\rho_0, \Omega)$-dominant projection if $\|\Pi_\Omega u_0\|_{\dot{H}^{1/2}}^2 \geq \rho_0 \|u_0\|_{\dot{H}^{1/2}}^2$.
\end{definition}

\begin{proposition}[Dominant Projection Dynamics]\label{prop:dominant}
If $u_0$ has $(\rho_0, \Omega)$-dominant projection with $\rho_0 > 1 - \varepsilon_*$, then:
\begin{enumerate}
\item The solution starts concentrated
\item Either concentration breaks (Case (ii)) or energy decreases (Case (iii))
\item Global regularity follows
\end{enumerate}
\end{proposition}

\begin{proof}
By assumption, $\rho(0) \geq \rho_0 > 1 - \varepsilon_*$, so the solution starts in the concentrated regime.

The Closing Lemma dichotomy applies:
\begin{itemize}
\item If $\lambda(0) < \lambda_{crit}$: concentration decreases ($d\rho/dt < 0$)
\item If $\lambda(0) \geq \lambda_{crit}$: energy decreases ($dE/dt < 0$)
\end{itemize}

In the first case, $\rho$ decreases until the solution exits concentration. In the second case, energy decreases until... it either exits concentration or decays to zero.

In all scenarios, blow-up is prevented.
\end{proof}

\subsection{Class 4: Solutions with Bounded Enstrophy Growth}

\begin{definition}
A solution has $K$-bounded enstrophy growth if $\|\omega(t)\|_{L^2} \leq K\|\omega(0)\|_{L^2}$ for all $t \in [0, T]$.
\end{definition}

\begin{proposition}
Solutions with bounded enstrophy growth cannot blow up (by BKM criterion), and automatically satisfy the corridor conditions with quantitative bounds.
\end{proposition}

\begin{proof}
BKM: blow-up requires $\int_0^T \|\omega\|_{L^\infty} dt = \infty$.

By Sobolev: $\|\omega\|_{L^\infty} \lesssim \|\omega\|_{H^{3/2}}$.

Bounded enstrophy $\|\omega\|_{L^2}$ plus bounded $\|u\|_{\dot{H}^{1/2}}$ (which follows from bounded enstrophy via interpolation) gives bounded $\|\omega\|_{H^{3/2}}$.

So $\int_0^T \|\omega\|_{L^\infty} dt < \infty$, and no blow-up.
\end{proof}

%==============================================================================
\section{Summary}
%==============================================================================

\begin{theorem}[Main Conclusions]
\mbox{}
\begin{enumerate}
\item \textbf{Forward stability:} The corridor condition cannot fail instantaneously. Solutions transition continuously between cases, with rate bounds given by Theorem~\ref{thm:short-time}.

\item \textbf{Blow-up contradiction:} Under hypothetical blow-up, the trilinear term $\mathcal{T}$ would need to satisfy $\mathcal{T} < -\nu D$, but the Closing Lemma shows this is incompatible with maintaining concentration. The ``escape routes'' (low $\lambda$ or high $\lambda$) both lead away from blow-up.

\item \textbf{Natural satisfaction:} The corridor condition automatically implies no blow-up for:
\begin{itemize}
\item Almost-helical flows (remain non-concentrated)
\item Spectrally localized data with high minimum frequency (energy decreases)
\item Data with dominant projected component (dichotomy applies)
\item Solutions with bounded enstrophy growth (BKM satisfied)
\end{itemize}
\end{enumerate}
\end{theorem}

\begin{remark}[What Remains]
The analysis shows that the corridor condition has good stability properties and correctly excludes blow-up for several natural subclasses. The main unresolved question is:

\textbf{Is concentration-compactness theory sufficient to guarantee that blow-up solutions concentrate?}

If yes, the proof is complete. If there are exotic blow-up scenarios that avoid concentration, additional analysis would be needed.

The standard concentration-compactness results (Kenig-Merle, Gallagher-Koch-Planchon) strongly suggest that concentration is unavoidable for critical blow-up, making the argument robust.
\end{remark}

\end{document}
